\begin{figure}[t]
    \centering
    \includegraphics[width=\linewidth]{iclr2026/figures/transfer_v13/transfer_map_17.pdf}
    \caption{\textbf{Generation-to-understanding transfer across visual capabilities.}
    Each column corresponds to an I2I supervision task and each row to an I2T capability in \methodname.
    Cells report the change in accuracy, in percentage points, relative to the I2T-only baseline.
    Blue indicates positive transfer and red indicates negative transfer.
    Outlined values mark each capability's largest observed gain, including ties.
    Stars mark $p<0.05$ from paired permutation tests against the I2T-only baseline.
    We show capabilities with more than 100 evaluation samples; Appendix~\ref{app:transfer} gives the full matrix and paired significance tests.}
    \label{fig:transfer-map}
\end{figure}

\section{Which generation data help which visual capabilities?}
\label{sec:transfer}
The controlled study shows that I2I training can improve the corresponding I2T task. We now evaluate the 15 I2I task types in \methodname and two additional objectives, inpainting and localization, across its 25 I2T capabilities.

\subsection{Experimental setup}
We follow the I2I $\rightarrow$ I2T recipe: starting from the same BAGEL checkpoint, the core experiments use approximately 50k I2I examples followed by 50k LLaVA-Instruct examples~\citep{liu2024llava}. We extend the map with inpainting and localization checkpoints; Appendix~\ref{app:transfer-setup} specifies the expanded checkpoint set. The baseline is trained only on the I2T data. We evaluate all checkpoints on the same 9,444 evaluation examples and define transfer from source $t$ to capability $c$ as
\[
\Delta(t,c)
=\operatorname{Acc}(M_t,\mathcal D_c)
-\operatorname{Acc}(M_{\mathrm{I2T}},\mathcal D_c),
\]
reported in percentage points. Figure~\ref{fig:transfer-map} shows the 19 capabilities with more than 100 evaluation examples; Appendix~\ref{app:transfer} gives all 25 capabilities, baseline accuracies, and benchmark-level results. Each checkpoint uses a single I2I source, separating the observed effects of different supervision objectives.

\subsection{Transfer depends on the source and target tasks}
\paragraph{Useful source-target pairs can share visual requirements.}
Counting improves with all 17 sources, with 12 gains reaching $p<0.05$. Object pointing produces the largest observed gain ($+3.57$ pp, $p<0.001$): both tasks require identifying the objects to be counted. Geometry supervision also benefits metric 3D relation: Z-depth, Euclidean depth, and surface normals improve accuracy by $2.90$, $2.77$, and $2.49$ pp, respectively (all $p<0.03$). These results connect shared visual requirements to useful transfer.

\paragraph{Different sources benefit complementary capabilities.}
Inpainting improves category recognition ($+3.22$ pp) and 2D ordering ($+11.05$ pp), while 2.5D segmentation produces the largest observed ordering gain ($+11.58$ pp). Colorization gives the largest observed gain on visual correspondence ($+5.71$ pp), crossing the taxonomy's Reconstruction--Reorganization boundary. All four gains have $p<0.01$. Useful generation supervision thus extends across recognition, geometry, and spatial reasoning, with different sources benefiting different capabilities.

All 17 checkpoints achieve higher pooled accuracy than the I2T-only baseline. Colorization and inpainting reach $70.08\%$ and $69.97\%$, respectively, compared with the baseline's $68.90\%$; both gains remain significant after Holm correction across the 17 sources (Appendix~\ref{app:transfer-tests}).

\begin{findingbox}
Image-generation supervision benefits complementary visual capabilities: useful transfer can follow shared visual requirements and cross task-family boundaries.
\end{findingbox}
